\chapter{Neural Networks and Backpropagation}

An intrusion prevention system observes four quantities on every connection: the number of
requests per second, the connection duration, the number of rejected packets, and the average
payload size. An analyst trying to separate malicious traffic from ordinary traffic finds a
pattern that is awkward to read. Connections with many requests but no rejected packets are
usually legitimate, because that is what fetching a large file looks like. Connections with many
rejected packets but few requests are also usually legitimate, because that is what a failed scan
looks like. The combination of the two is almost always an attack. Each feature by itself
separates nothing, whereas their combination separates sharply.

A linear threshold cannot draw such a separator, and the first part of this chapter shows that
the obstacle is not the amount of data. What must change is the shape of the function. One hidden
layer is already enough for that pattern. Once the shape has changed, the next question appears,
and it is more serious: how do we tune the thirteen parameters of that small network without
deriving a new formula for every architecture?

The answer is the chain rule of Chapter 12, run backward on the computational graph. This chapter
carries that rule down to the numbers: one forward step with given weights, all thirteen
gradients computed by hand and then checked with NumPy, one gradient-descent step, and the
gradient with respect to the input that later becomes a weapon in Chapter 29. What this chapter
promises is not a clever network, but a differentiation engine that can be checked line by line.

\begin{untukanak}
Imagine an assembly line with several workers. The first worker works on the raw material, the
second works on the output of the first, and so on until the last worker hands the finished item
to an inspector. The inspector weighs the item, then says: too heavy, two hundred grams over.

Now the question: who has to change the way they work? The first worker? The last worker? If the
foreman can change only one setting and then rerun the whole line, he needs two tries per setting.
With thirteen settings, that is twenty-six runs of the line for a single fix.

It turns out there is a far cleverer way. After one run, we can walk back from the inspector to
the last worker, then to the previous worker, all the way to the first. At each stop we compute
one number: how much the final error changes if this person changes their output a little. With a
single backward pass, every worker knows their own share, and nothing has to be run again.

In this chapter we will learn why information can flow backward even though the goods only flow
forward, and why that backward pass gives an answer exactly as accurate as running the line
twenty-six times.
\end{untukanak}

\section{Understanding bridge: from story to symbol}
\tingkat{basic}

The formulas in this chapter are easier to read if the order is reversed: a story without
symbols, then arithmetic on paper, then symbols as shorthand for steps we have already
performed ourselves.

\subsection{Plain analogy: how much, not who}

The assembly-line analogy answered ``who should change''. It did not answer ``how much''.
Picture three taps in a row above one bucket, with a target line on the wall. The third tap
feels closest, but it can only pass on what the second tap sends; the first tap matters
most, yet its effect depends on the two taps between it and the bucket. What we need is not
turning every tap and re-measuring, but one number per tap: how many millimetres does the
water rise if this tap moves a little? One pass, traced backwards from the bucket.

\subsection{Compute by hand first}

\begin{contoh}
One neuron, four numbers: weight $w = 0.5$, bias $b = 0.1$, input $x = 2.0$, correct
answer $y = 1$, learning rate $\eta = 0.1$. Follow the steps in order; every number here
can be checked on an ordinary calculator.

\begin{enumerate}
  \item Raw score: $w x = 0.5 \times 2.0 = 1.0$, plus the bias gives $z = 1.1$.
  \item Probability: $e^{-1.1} = 0.33287$, so $\sigma = 1/(1+0.33287) = 0.7502601$.
  \item Loss: because $y = 1$ only one term survives, $L = -\ln 0.7502601 = 0.2873353$.
  \item Gap between probability and truth: $\sigma - y = 0.7502601 - 1 = -0.2497399$.
  \item Chain rule to the weight: multiply that gap by the input, $-0.2497399 \times 2.0
        = -0.4994798$. The bias is not multiplied by anything, so it keeps $-0.2497399$.
  \item Move: $w' = 0.5 - 0.1(-0.4994798) = 0.5499480$ and $b' = 0.1 - 0.1(-0.2497399)
        = 0.1249740$. The new score is $z' = 1.2248699$, the new probability $0.7729194$,
        the new loss $0.2575805$: down $0.0297549$ from $0.2873353$.
\end{enumerate}

Mind the sign in step 5: the gradient is negative, and because we \emph{subtract} it, the
weight rises. The minus sign in the update rule is not decoration; it reverses the
direction.
\end{contoh}

One check you must not skip: replace the analytic derivative by a finite difference, that
is, evaluate the loss at $w + 0.000001$ and $w - 0.000001$ and divide the difference by
$0.000002$. The answer is $-0.4994798$, equal to seven digits. When the hand result and the
numeric result disagree, the sign is almost always what is wrong.

\subsection{Symbol dictionary}

The symbols below appear in any book on neural networks. Nothing here needs memorising:
reading each one aloud takes a single sentence, and that sentence is what we computed above.

\begin{description}
  \item[$\partial L/\partial w$] read ``del L del w'': if weight $w$ moves a little, how
        much does the loss $L$ change? In our example the answer is $-0.4994798$ (step 5).
  \item[$z = Wx + b$] read ``z equals W times x plus b'': the raw score before it is
        squeezed into a probability. The one-number version of this formula is step 1.
  \item[$\sigma(z)$] read ``sigma of z'': the S-shaped function that keeps the score
        between zero and one. In step 2 it gives $0.7502601$.
  \item[$\nabla_w L$] read ``nabla w L'': every derivative of the loss with respect to
        every weight at once; its arrow answers ``which way does the loss fall fastest''.
        Its entries are numbers like the one in step 5, only many.
\end{description}

\subsection{Why does this matter for security?}

The same derivative that repairs a model also breaks it. An attacker needs one number: if
one input feature moves a little, how much does the loss rise? In our example $x$ has
derivative $-0.1248699$, so shrinking $x$ raises the loss. Moving it by $1.9$ sends $x$ to
$0.1$, the probability drops to $0.5374298$, and the loss rises from $0.2873353$ to
$0.6209570$, a factor of $2.16$. One feature, one direction, one step.

For the defender the same question runs backwards: which feature has the largest
derivative? That feature is the cheapest to fake, so it most needs rounding, clipping, or
sanity checks. Chapter 29 uses this recipe in full, with perturbation size $\epsilon$ and
the formula $x_{\mathrm{adv}} = x + \epsilon \cdot \operatorname{sign}(\nabla_x L)$.

\begin{catatan}
A warning that is easily missed: on real systems an attacker rarely computes the gradient
exactly; it is estimated by changing the input many times. That estimate is more expensive
and noisier, but it still works. A closed model is therefore not a defence.
\end{catatan}

\subsection{Hand computation drills}

Work with a calculator, write on paper, then compare. Answers sit in parentheses.

\begin{catatan}
\begin{enumerate}
  \item Repeat the whole example with $b = 0.3$, everything else unchanged. Compute $z$,
        $\sigma$, $L$, $\partial L/\partial w$, $\partial L/\partial b$.
        ($z = 1.3$; $\sigma = 0.7858350$; $L = 0.2410085$; $-0.4283300$; $-0.2141650$.)
  \item Repeat the original example with $\eta = 0.5$. What are $w'$, $b'$, and the new
        loss? ($w' = 0.7497399$; $b' = 0.2248699$; $L' = 0.1640633$, down $0.1232720$.)
  \item Repeat the original example with $y = 0$, everything else unchanged. What is
        $\partial L/\partial w$ now, and which way does the weight move?
        ($0.7502601$ and $1.5005202$; positive, so the weight falls.)
  \item In the original example, move $x$ by $\epsilon = 0.5$ in the direction that raises
        the loss. What are the new $x$ and the new loss? ($x = 1.5$; $L = 0.3558651$.)
\end{enumerate}
\end{catatan}

\subsection{From paper to code}

The code below runs the same six steps, with variable names that point at their step.
Nothing here is new: every line is one of the lines above.

\begin{lab}
\begin{verbatim}
import numpy as np

# One example, the same numbers as in "Compute by hand first".
x, y = 2.0, 1.0
w, b, eta = 0.5, 0.1, 0.1

# Step 1: z = w*x + b  (raw score)
z = w*x + b
print("Step 1  z        =", round(z, 7))

# Step 2: sigma(z) = 1/(1+e^-z)  (squeeze the score into 0..1)
s = 1/(1+np.exp(-z))
print("Step 2  sigma    =", round(s, 7))

# Step 3: L = -ln(s)  (loss; y=1 leaves a single term)
L = -np.log(s)
print("Step 3  L        =", round(L, 7))

# Step 4: dL/dz = sigma - y  (probability minus reality)
dLdz = s - y
print("Step 4  dL/dz    =", round(dLdz, 7))

# Step 5: chain rule down to weight and bias
dLdw, dLdb = dLdz*x, dLdz
print("Step 5  dL/dw    =", round(dLdw, 7), " dL/db =", round(dLdb, 7))

# Step 6: one gradient-descent step
w2, b2 = w - eta*dLdw, b - eta*dLdb
s2 = 1/(1+np.exp(-(w2*x + b2)))
print("Step 6  w', b'   =", round(w2, 7), round(b2, 7),
      " L' =", round(-np.log(s2), 7))

# Check: analytic gradient against a finite difference.
def loss(ww, bb):
    return -np.log(1/(1+np.exp(-(ww*x + bb))))
h = 1e-6
print("finite difference dL/dw =", round((loss(w+h, b)-loss(w-h, b))/(2*h), 7),
      " dL/db =", round((loss(w, b+h)-loss(w, b-h))/(2*h), 7))

# Gradient with respect to the INPUT: this is what attacks a model (Chapter 29).
dLdx = dLdz*w
xa = x + 1.9*np.sign(dLdx)          # the sign of the gradient sets the direction
sa = 1/(1+np.exp(-(w*xa + b)))
print("dL/dx =", round(dLdx, 7), " x_adv =", round(xa, 4),
      " sigma =", round(sa, 7), " L =", round(-np.log(sa), 7))
\end{verbatim}

The output must match the paper: $1.1$, $0.7502601$, $0.2873353$, $-0.2497399$,
$-0.4994798$, then $0.2575805$. If one number is off, do not patch the code: find the step
that differs on paper.
\end{lab}

\section{Linear thresholds and the perceptron}
\tingkat{basic}

The oldest form of classifier is the linear threshold. For an input
$\mathbf{x}\in\mathbb{R}^{d}$, weights $\mathbf{w}\in\mathbb{R}^{d}$, and bias
$b\in\mathbb{R}$, we compute a single score and decide the class from the sign of that score.

\kamus{perceptron}{perceptron}{A classifier that decides only from which side of a single
separating line or plane a point falls.}

\begin{definisi}[Linear threshold]
A function $f:\mathbb{R}^{d}\to\{-1,+1\}$ is called a linear threshold if there are
$\mathbf{w}$ and $b$ such that $f(\mathbf{x})=\operatorname{sign}(\mathbf{w}^{\top}\mathbf{x}+b)$.
The set $\{\mathbf{x}:\mathbf{w}^{\top}\mathbf{x}+b=0\}$ is the separating hyperplane, and that
hyperplane is the only form of separator this model can learn.\index{perceptron}
\end{definisi}

Perceptron training works by correcting the weights only when there is a mistake. For a label
$y\in\{-1,+1\}$ and an example $\mathbf{x}$, if $y\,\mathbf{w}^{\top}\mathbf{x}\le 0$ then we
update
\[
  \mathbf{w}\leftarrow\mathbf{w}+y\,\mathbf{x}.
\]
Note that there is no learning rate here: the size of the step is determined by the size of the
example, and training stops by itself as soon as all examples are correct. Does it always stop?

\begin{teorema}[Perceptron convergence]
Suppose every example satisfies $\lVert\mathbf{x}_{i}\rVert\le1$ and there is a unit vector
$\mathbf{w}^{*}$ and $\gamma>0$ such that $y_{i}\,\mathbf{w}^{*\top}\mathbf{x}_{i}\ge\gamma$ for
every $i$. Then the perceptron algorithm starting from $\mathbf{w}_{0}=\mathbf{0}$ makes at most
$\lfloor1/\gamma^{2}\rfloor$ updates.\index{perceptron}
\end{teorema}

\begin{proof}
Suppose an update occurs $k$ times, on the examples $\mathbf{x}_{1},\dots,\mathbf{x}_{k}$. Every
update occurs only when that example is misclassified, so
$y_{j}\,\mathbf{w}_{j-1}^{\top}\mathbf{x}_{j}\le0$. Hence
\[
  \lVert\mathbf{w}_{j}\rVert^{2}
  =\lVert\mathbf{w}_{j-1}\rVert^{2}+2y_{j}\mathbf{w}_{j-1}^{\top}\mathbf{x}_{j}
   +\lVert\mathbf{x}_{j}\rVert^{2}
  \le\lVert\mathbf{w}_{j-1}\rVert^{2}+1 .
\]
By induction, $\lVert\mathbf{w}_{k}\rVert^{2}\le k$.

On the other side, the inner product with $\mathbf{w}^{*}$ grows by at least $\gamma$ each
time:
\[
  \mathbf{w}^{*\top}\mathbf{w}_{j}
  =\mathbf{w}^{*\top}\mathbf{w}_{j-1}+y_{j}\mathbf{w}^{*\top}\mathbf{x}_{j}
  \ge\mathbf{w}^{*\top}\mathbf{w}_{j-1}+\gamma ,
\]
so $\mathbf{w}^{*\top}\mathbf{w}_{k}\ge k\gamma$. Since $\mathbf{w}^{*}$ is a unit vector, the
Cauchy-Schwarz inequality gives
\[
  k\gamma\le\mathbf{w}^{*\top}\mathbf{w}_{k}\le\lVert\mathbf{w}_{k}\rVert\le\sqrt{k}.
\]
After dividing by $\sqrt{k}$ and squaring, $k\gamma^{2}\le1$, that is $k\le1/\gamma^{2}$.
\end{proof}

That bound does not depend on the dimension or on the number of examples, and that is its
appeal. But when $\gamma$ does not exist, this theorem says nothing, and the perceptron really
does oscillate forever; it is also not a statement about generalization.

Now a more fundamental bound. Four points in the plane are already enough to show that not every
labeling can be separated by a single hyperplane. The points $(0,0)$ and $(1,1)$ have label $0$,
the points $(1,0)$ and $(0,1)$ have label $1$. This pattern is usually called XOR.

\begin{proposisi}
The four XOR points are not linearly separable.
\end{proposisi}

\begin{proof}
Suppose there are $w_{1},w_{2},b$ that separate them, with label $1$ on the positive side. Label
$0$ at $(0,0)$ gives $b\le0$. Label $1$ at $(1,0)$ gives $w_{1}+b>0$, hence $w_{1}>-b\ge0$. By
the same route, $w_{2}>-b\ge0$. Therefore $w_{1}+w_{2}>-2b$, and
\[
  w_{1}+w_{2}+b>-b\ge0 .
\]
But label $0$ at $(1,1)$ demands $w_{1}+w_{2}+b\le0$. Contradiction.
\end{proof}

This contradiction holds for every choice of $\mathbf{w}$ and $b$, not only for some particular
training outcome. Adding examples, changing the learning rate, or normalizing the features will
not help. What must change is the shape of the function, and the smallest change available is to
insert one layer in between.

\kamus{hidden layer}{lapisan tersembunyi}{A layer of units that the outside world does not see:
its input comes from the previous layer, its output goes to the next layer.}

A hidden layer works in two steps. First, compute
$\mathbf{z}_{1}=\mathbf{W}_{1}\mathbf{x}+\mathbf{b}_{1}$ with
$\mathbf{W}_{1}\in\mathbb{R}^{2\times4}$, then apply
$\operatorname{ReLU}(t)=\max\{t,0\}$ to every component, and write the result $\mathbf{h}$.
Second, compute a linear score from $\mathbf{h}$. For XOR, with
\[
  \mathbf{W}_{1}=\begin{pmatrix}1&1\\1&1\end{pmatrix},\qquad
  \mathbf{b}_{1}=\begin{pmatrix}0\\-1\end{pmatrix},\qquad
  \mathbf{w}_{2}=\begin{pmatrix}1\\-2\end{pmatrix},\qquad b_{2}=0,
\]
the four points give: $(0,0)$ produces $\mathbf{z}_{1}=(0,-1)$ and $\mathbf{h}=(0,0)$ with score
$0$; the points $(1,0)$ and $(0,1)$ both produce $\mathbf{z}_{1}=(1,0)$ and $\mathbf{h}=(1,0)$
with score $1$; the point $(1,1)$ produces $\mathbf{z}_{1}=(2,1)$ and $\mathbf{h}=(2,1)$ with
score $1(2)-2(1)=0$. With a threshold of $0.5$ all four points are classified correctly.

What makes this work is not magic but simple arithmetic. The first unit computes
$\max\{x_{1}+x_{2},0\}$ and the second computes $\max\{x_{1}+x_{2}-1,0\}$. Both are linear, but
their difference is a piecewise function that turns at two places. A linear combination of linear
functions is still linear; a linear combination of piecewise functions is not.

Such a boundary is made of line segments. Each unit in the first layer contributes one straight
line, and the second layer chooses which side is used, so the final boundary is a polygon;
Figure~\ref{fig:dua-lapis} shows its shape.

\begin{figure}[htbp]
\centering
\begin{tikzpicture}[>=latex, font=\small]
  % left panel: two hidden layers
  \foreach \i/\yy in {1/1.5,2/-0.5}
    \node[circle,draw=layergrey,fill=layergrey!10,inner sep=1.7pt] (I\i) at (0.4,\yy) {};
  \foreach \i/\yy in {1/2.4,2/1.65,3/0.9,4/0.15,5/-0.6}
    \node[circle,draw=layerteal,fill=layerteal!12,inner sep=1.7pt] (H1\i) at (2.1,\yy) {};
  \foreach \i/\yy in {1/1.6,2/0.7,3/-0.2}
    \node[circle,draw=layerteal,fill=layerteal!12,inner sep=1.7pt] (H2\i) at (3.7,\yy) {};
  \node[circle,draw=layerindigo,fill=layerindigo!12,inner sep=1.7pt] (O) at (5.1,0.7) {};
  \foreach \i in {1,2} \foreach \j in {1,...,5} \draw[layergrey!45] (I\i) -- (H1\j);
  \foreach \i in {1,...,5} \foreach \j in {1,2,3} \draw[layergrey!45] (H1\i) -- (H2\j);
  \foreach \j in {1,2,3} \draw[layergrey!45] (H2\j) -- (O);
  \node[font=\scriptsize] at (0.4,-1.35) {input};
  \node[font=\scriptsize] at (2.1,-1.35) {five units};
  \node[font=\scriptsize] at (3.7,-1.35) {three units};
  \node[font=\scriptsize] at (5.1,-1.35) {output};
  % panel kanan: bidang input
  \draw[layergrey] (6.7,-0.9) rectangle (14.4,2.9);
  \node[font=\scriptsize] at (14.1,-1.35) {$x_{1}$};
  \node[font=\scriptsize] at (6.4,2.75) {$x_{2}$};
  \fill[layerorange!18] (8.3,0.3) -- (9.7,2.4) -- (12.6,2.55) -- (13.3,1.1) -- (11.2,0.0) -- cycle;
  \draw[layerorange,thick] (8.3,0.3) -- (9.7,2.4) -- (12.6,2.55) -- (13.3,1.1) -- (11.2,0.0) -- cycle;
  \draw[dashed,layergrey!70] (7.50,-0.9) -- (11.43,2.9);
  \draw[dashed,layergrey!70] (6.70,2.23) -- (14.40,2.64);
  \draw[dashed,layergrey!70] (12.43,2.9) -- (14.27,-0.9);
  \draw[dashed,layergrey!70] (14.40,1.68) -- (9.48,-0.9);
  \draw[dashed,layergrey!70] (6.70,0.50) -- (14.40,-0.33);
  \foreach \xx/\yy in {9.9/1.6, 11.4/1.3, 10.6/0.6}
    \draw[accent,thick] (\xx-0.09,\yy-0.09) -- (\xx+0.09,\yy+0.09)
      (\xx-0.09,\yy+0.09) -- (\xx+0.09,\yy-0.09);
  \foreach \xx/\yy in {7.7/1.4, 8.5/-0.3, 12.3/-0.5, 13.6/2.2}
    \draw[layergrey,fill=white] (\xx,\yy) circle (2.1pt);
  \draw[accent,thick] (11.0,-1.30) -- (11.15,-1.15) (11.0,-1.15) -- (11.15,-1.30);
  \node[font=\scriptsize,anchor=west] at (11.2,-1.23) {class 1};
  \draw[layergrey,fill=white] (13.0,-1.23) circle (2.1pt);
  \node[font=\scriptsize,anchor=west] at (13.2,-1.23) {class 0};
\end{tikzpicture}
\caption{A sketch of how two hidden layers turn one line into a fold. Left: each unit in the
first layer pulls one straight line across the input plane, and the second layer combines the
sides of those lines. Right: the decision boundary is a five-sided polygon, drawn thick; the
dashed lines are the extension of each side, namely the line belonging to one unit. Crosses are
class 1 inside the polygon, circles class 0 outside. The figure is a sketch, not a training
result: the number of sides follows the number of units in the first layer.}
\label{fig:dua-lapis}
\end{figure}


\section{Activation functions and two diseases of the inner layers}
\tingkat{intermediate}

Without a nonlinear function between the layers, the whole network collapses into a single
matrix. If ReLU is replaced by the identity, the output is
$\mathbf{W}_{2}\mathbf{W}_{1}\mathbf{x}$, that is, a single matrix product
$\mathbf{W}\mathbf{x}$ with $\mathbf{W}=\mathbf{W}_{2}\mathbf{W}_{1}$. Every hidden layer
therefore needs a nonlinear activation, and the choice of activation determines which disease
will appear when the network is made deep.

\kamus{activation function}{fungsi aktivasi}{The gate that determines how much of a unit's signal
is passed on, and at the same time the one that fixes the shape of its derivative.}

\begin{definisi}[ReLU, sigmoid, and tanh]
For $t\in\mathbb{R}$,
\[
  \operatorname{ReLU}(t)=\max\{t,0\},\qquad
  \sigma(t)=\frac{1}{1+e^{-t}},\qquad
  \tanh(t)=\frac{e^{t}-e^{-t}}{e^{t}+e^{-t}} .
\]
The derivative of ReLU is $1$ for $t>0$ and $0$ for $t<0$; at $t=0$ the derivative does not
exist, and this book uses the value $0$ as a convention. The derivative of tanh satisfies
$\tanh'(t)=1-\tanh^{2}(t)$.\index{activation function}\index{relu}
\end{definisi}

The value $0$ at the break point of ReLU is a convention: the set where the derivative fails is
only $\{0\}$, of measure zero, so it does not affect the derivative almost everywhere. The
practical consequence: a unit whose input is always negative receives zero gradient and will
never change.

Sigmoid has a neater property, and that property is the source of the trouble.

\begin{lemma}[Derivative of sigmoid]
For every $t$ we have $\sigma'(t)=\sigma(t)\bigl(1-\sigma(t)\bigr)$, hence
$0<\sigma'(t)\le\tfrac14$.\index{saturation}
\end{lemma}

\begin{proof}
Write $\sigma(t)=(1+e^{-t})^{-1}$. Then
\[
  \sigma'(t)=-(1+e^{-t})^{-2}\cdot(-e^{-t})
  =\frac{e^{-t}}{(1+e^{-t})^{2}}
  =\frac{1}{1+e^{-t}}\cdot\frac{e^{-t}}{1+e^{-t}}
  =\sigma(t)\bigl(1-\sigma(t)\bigr).
\]
Since $0<\sigma(t)<1$, the product is positive. Writing $s=\sigma(t)$ gives
$s(1-s)=\tfrac14-(s-\tfrac12)^{2}\le\tfrac14$, and equality is attained when $s=\tfrac12$.
\end{proof}

The bound $\tfrac14$ is the one to hold on to. The following table computes $\sigma$,
$\sigma'$, and $\tanh$ at several values; the numbers were produced by NumPy and are written to
six decimals.

\begin{center}
\small
\begin{tabular}{rrrrr}
\hline
$t$ & $\sigma(t)$ & $\sigma'(t)$ & $\tanh(t)$ & $1-\tanh^{2}(t)$ \\
\hline
$0$  & $0.500000$ & $0.250000$ & $0.000000$ & $1.000000$ \\
$1$  & $0.731059$ & $0.196612$ & $0.761594$ & $0.419974$ \\
$2$  & $0.880797$ & $0.104994$ & $0.964028$ & $0.070650$ \\
$5$  & $0.993307$ & $0.006648$ & $0.999909$ & $0.000181$ \\
$10$ & $0.999955$ & $0.000045$ & $0.999999$ & $0.000002$ \\
\hline
\end{tabular}
\end{center}

The fourth and fifth columns show what is called saturation: once $t$ moves away from zero, the
output barely changes with the input, and the derivative tends to zero. At $t=10$ the derivative
is only $4.54\times10^{-5}$, about one twenty-thousandth of its maximum.

The consequence can be computed. If the error propagates backward through $L$ consecutive layers
whose derivative is at most $\tfrac14$, the remaining signal is at most $4^{-L}$: for $L=5$ about
$9.77\times10^{-4}$, and for $L=10$ about $9.54\times10^{-7}$. A layer far from the output
receives a signal that is only one millionth, so it practically does not learn. This symptom is
called the vanishing gradient. ReLU avoids it for a different reason: its derivative is $1$ on
the positive side, so there is no shrinkage, but its derivative is $0$ on the negative side, so
there is unit death. Both diseases are not implementation errors; both are consequences of the
shape of the function.\index{vanishing gradient}

One more statement is often quoted beyond its content: one hidden layer is enough for all
continuous functions.

\begin{teorema}[Universal approximation, statement]
Suppose $\varphi$ is bounded, monotonically increasing, and has different limits at $-\infty$
and $+\infty$. For every compact set $K\subset\mathbb{R}^{d}$, every continuous function
$f:K\to\mathbb{R}$, and every $\varepsilon>0$, there are $N$ and parameters
$\mathbf{w}_{i},b_{i},v_{i}$ such that
\[
  \sup_{\mathbf{x}\in K}\Bigl|f(\mathbf{x})-\sum_{i=1}^{N}v_{i}\,
  \varphi(\mathbf{w}_{i}^{\top}\mathbf{x}+b_{i})\Bigr|<\varepsilon .
\]
\end{teorema}

The proof of this statement is not included here; the original proof is in Cybenko (1989) and a
more general version in Hornik (1991), and the extension to ReLU is in Leshno et al. (1993).
What is important to read carefully is what the theorem does \emph{not} say. It states that such
a set of functions is dense, not that a particular network equals $f$. It does not say how large
$N$ must be for error $\varepsilon$ in dimension $d$; in the worst case $N$ grows exponentially
with $d$ (Goodfellow et al., section 6.4.1). It also does not say that those parameters can be
found by gradient descent, and it says nothing about data not yet seen. The sentence ``a network
can approximate anything'' is therefore true as a statement of density and misleading as a
practical promise.

\section{From score to loss}
\tingkat{intermediate}

The network's output must be compared with the label through a single number that can be
minimized, and the choice of that number determines how much signal reaches the lower layers.
For binary classification with label $y\in\{0,1\}$ and output $a=\sigma(z_{2})$, the loss used
is the binary cross-entropy
\[
  L(a,y)=-\bigl[y\ln a+(1-y)\ln(1-a)\bigr].
\]
Cross-entropy is the negative logarithm of the confidence the model gives to the correct class,
so a confident mistake is punished far more heavily than a hesitant one.

\kamus{cross-entropy}{entropi silang}{A measure of surprise: the more confident the model is in
a wrong answer, the larger the penalty, and the penalty grows without bound.}\index{cross-entropy}

\begin{proposisi}[Gradient of cross-entropy]
For $a=\sigma(z)$ and $L$ the binary cross-entropy we have
$\dfrac{\partial L}{\partial z}=a-y$.
\end{proposisi}

\begin{proof}
Differentiate $L$ with respect to $a$:
\[
  \frac{\partial L}{\partial a}=-\frac{y}{a}+\frac{1-y}{1-a}
  =\frac{a-y}{a(1-a)} .
\]
By the chain rule and the lemma on the derivative of the sigmoid,
$\dfrac{\partial a}{\partial z}=a(1-a)$, so
\[
  \frac{\partial L}{\partial z}=\frac{a-y}{a(1-a)}\cdot a(1-a)=a-y .
\]
Numerator and denominator cancel, and that is the point.
\end{proof}

Compare with the squared loss $L=\tfrac12(a-y)^{2}$, which gives
$\partial L/\partial z=(a-y)\,a(1-a)$. The factor $a(1-a)$ does not cancel; it shrinks exactly
when the model is most wrong with most confidence. For $z=10$ and $y=0$,
$a(1-a)=4.5396\times10^{-5}$, so the signal to the weights is only $4.5394\times10^{-5}$ even
though the error is almost one. With cross-entropy the signal stays $a-y=0.999955$. The ratio of
the two numbers is about $2.2\times10^{4}$; this is why the sigmoid output is paired with
cross-entropy, not with the square.

For $K$ classes, the same pairing holds, and this time we prove it.

\kamus{softmax}{softmax}{Turns a list of scores into a division of probabilities that sums to one,
by comparing every score with the sum of them all.}\index{softmax}

\begin{proposisi}[Softmax and cross-entropy]
Suppose $\mathbf{p}=\operatorname{softmax}(\mathbf{z})$ with
$p_{k}=e^{z_{k}}/\sum_{j}e^{z_{j}}$, and $L=-\sum_{k}y_{k}\ln p_{k}$ with
$\sum_{k}y_{k}=1$. Then $\dfrac{\partial L}{\partial z_{i}}=p_{i}-y_{i}$.
\end{proposisi}

\begin{proof}
The component derivative of softmax is
$\dfrac{\partial p_{k}}{\partial z_{i}}=p_{k}(\delta_{ki}-p_{i})$, with $\delta_{ki}$ the
Kronecker delta. Hence
\[
  \frac{\partial L}{\partial z_{i}}
  =-\sum_{k}\frac{y_{k}}{p_{k}}\cdot p_{k}(\delta_{ki}-p_{i})
  =-\sum_{k}y_{k}\delta_{ki}+p_{i}\sum_{k}y_{k}
  =-y_{i}+p_{i},
\]
because $\sum_{k}y_{k}=1$.
\end{proof}

For two classes, $\operatorname{softmax}((t,0))_{1}=\sigma(t)$, so this formula returns to $a-y$
as before. Note also that adding a constant $c$ to all $z_{k}$ does not change $\mathbf{p}$,
because the factor $e^{c}$ cancels; only differences of scores are meaningful, which is also why
adding a large constant to all logits damages floating-point accuracy (Chapter 22).

\begin{untukanak}
Have you ever played a number-guessing game where, every time you are wrong, you are only told
``too big'' or ``too small'', without being told by how much? If your guess is nine hundred and
the answer is one hundred, the remark ``too big'' is exactly the same as if your guess were one
hundred and one.

A neural network has the same problem, but more complicated. It has thirteen knobs, and every
knob affects the answer through several intermediaries. If we only know that the answer is
wrong, we do not know which knob to turn, in which direction, or by how much.

Fortunately, a wrong answer has a degree: how far off, and in which direction. With one measure
called the loss, we can ask every knob: if you were turned a little, does the loss go up or down,
and how fast? The answer is one number per knob.

In this chapter we will learn why a single loss measure is enough to give thirteen answers at
once, and why those answers can be trusted.
\end{untukanak}

\section{Backpropagation}
\tingkat{advanced}

A network is a composition of functions, so its derivative is computed with the chain rule of
Chapter 12. What makes it cheap is the order of work: from the output backward to the input, not
from the input forward to the output. For every intermediate quantity $u$ in the computational
graph, we assign one number stating how much the loss changes if $u$ changes.

\kamus{adjoint}{adjoint}{The image of how much the final error is affected by one quantity in the
middle of the network; one number per quantity.}

\begin{definisi}[Adjoint]
Suppose $C$ is the loss on one example and $u$ one of the intermediate quantities. The adjoint of
$u$ is $\bar{u}=\partial C/\partial u$.\index{adjoint}
\end{definisi}

\begin{proposisi}[Adjoint recursion]
If $u$ affects $C$ only through $v_{1},\dots,v_{m}$, then
$\bar{u}=\sum_{i=1}^{m}\bar{v}_{i}\,\dfrac{\partial v_{i}}{\partial u}$.
\end{proposisi}

\begin{proof}
This is the chain rule for functions of several variables (Chapter 12) applied to $C$ as a
function of $v_{1},\dots,v_{m}$, with each $v_{i}$ as a function of $u$. Since $u$ does not
affect $C$ in any other way, all other terms in that chain rule are zero.
\end{proof}

That recursion turns the work into a single pass. Starting from $\bar{C}=1$, we visit the
intermediate quantities in the reverse of the forward order. Every quantity receives
contributions from all the quantities that use it, so nothing is missed or counted twice. This
is backpropagation, and it is not a new algorithm: it is only the chain rule run on the graph in
the right order.\index{backpropagation}

Figure~\ref{fig:graf-balik} places both directions on one graph. Full arrows carry the forward
values; dashed arrows carry the adjoints, and the numbers on those arrows are taken from the
worked example in the next section.

\begin{figure}[htbp]
\centering
\begin{tikzpicture}[>=latex, font=\small]
  \node[draw=layerteal,rounded corners=2pt,fill=layerteal!10,inner sep=4pt] (X)  at (0,0)    {$\mathbf{x}$};
  \node[draw=layerteal,rounded corners=2pt,fill=layerteal!10,inner sep=4pt] (Z1) at (2.7,0)  {$\mathbf{z}_{1}$};
  \node[draw=layerteal,rounded corners=2pt,fill=layerteal!10,inner sep=4pt] (H)  at (5.3,0)  {$\mathbf{h}$};
  \node[draw=layerteal,rounded corners=2pt,fill=layerteal!10,inner sep=4pt] (Z2) at (7.5,0)  {$z_{2}$};
  \node[draw=layerteal,rounded corners=2pt,fill=layerteal!10,inner sep=4pt] (A)  at (9.3,0)  {$a$};
  \node[draw=layerteal,rounded corners=2pt,fill=layerteal!10,inner sep=4pt] (L)  at (11.4,0) {$L$};
  \node (W1) at (2.7,1.45) {$\mathbf{W}_{1},\mathbf{b}_{1}$};
  \node (W2) at (7.5,1.45) {$\mathbf{w}_{2},b_{2}$};
  \foreach \a/\b in {X/Z1,Z1/H,H/Z2,Z2/A,A/L} \draw[->,thick,ink] (\a) -- (\b);
  \draw[->,thick,ink] (W1) -- (Z1);
  \draw[->,thick,ink] (W2) -- (Z2);
  \draw[->,dashed,accent] (L.south)  .. controls +(0,-1.0) and +(0,-1.0) .. (Z2.south);
  \draw[->,dashed,accent] (Z2.south) .. controls +(0,-1.0) and +(0,-1.0) .. (H.south);
  \draw[->,dashed,accent] (H.south)  .. controls +(0,-1.0) and +(0,-1.0) .. (Z1.south);
  \draw[->,dashed,accent] (Z1.south) .. controls +(0,-1.0) and +(0,-1.0) .. (X.south);
  \draw[->,dashed,accent] (3.05,0.25) -- (3.05,1.25);
  \draw[->,dashed,accent] (7.85,0.25) -- (7.85,1.25);
  \node[font=\scriptsize] at (9.9,-1.4)  {$\bar{z}_{2}=a-y=-0.2554$};
  \node[font=\scriptsize] at (6.5,-2.0)  {$\bar{\mathbf{h}}=(-0.2043;\,0.1532)$};
  \node[font=\scriptsize] at (4.0,-1.4)  {$\bar{\mathbf{z}}_{1}=\bar{\mathbf{h}}$};
  \node[font=\scriptsize,align=center] at (1.1,-1.95)
    {$\nabla_{\mathbf{x}}L=(-0.0868;\,0.1328;$\\$(-0.0971;\,-0.1430)$};
  \node[font=\scriptsize] at (2.7,2.15) {$\nabla_{\mathbf{W}_{1}}L=\bar{\mathbf{z}}_{1}\mathbf{x}^{\top}$};
  \node[font=\scriptsize] at (9.0,2.15) {$\nabla_{\mathbf{w}_{2}}L=\bar{z}_{2}\mathbf{h}$};
\end{tikzpicture}
\caption{The computational graph of the 4--2--1 network on the worked example. Thick full arrows
are the forward direction: input and weights become $\mathbf{z}_{1}$, the ReLU gate becomes
$\mathbf{h}$, then the score $z_{2}$, the output $a$, and the loss $L$. Dashed arrows are the
backward direction that carries the adjoints: $\bar{z}_{2}=a-y=-0.2554$ enters the output
weights, $\bar{\mathbf{h}}=(-0.2043;\,0.1532)$ enters the gate, and
$\nabla_{\mathbf{x}}L=(-0.0868;\,0.1328;\,-0.0971;\,-0.1430)$ exits at the input. One backward
pass produces all thirteen gradients at once.}
\label{fig:graf-balik}
\end{figure}

\begin{bacakeras}
Two voices. First voice: ``The loss is too large.'' Second voice: ``Who is at fault?'' First
voice: ``Not one person. The fault is at the end, and the fault is divided backward, little by
little.''

Listen once more, slowly. The error is born at the end, in the last number. From there it walks
backward, stops at every junction, and asks: through here, how much is your share? The one
closest to the end answers largest. The one far behind answers smallest, because the answer has
already been filtered twice on the way home.

Second voice: ``So the ones in front are not at fault?'' First voice: ``They are at fault to the
extent of the share that reaches them. No more. That is why we do not accuse, we compute.''

Repeat three times until it sounds right: fault at the end, divided backward, little by little,
along the path.
\end{bacakeras}


Now derive it in full for the 4--2--1 network with ReLU in the hidden layer and sigmoid at the
output. The forward computation is
\[
  \mathbf{z}_{1}=\mathbf{W}_{1}\mathbf{x}+\mathbf{b}_{1},\qquad
  \mathbf{h}=\operatorname{ReLU}(\mathbf{z}_{1}),\qquad
  z_{2}=\mathbf{w}_{2}^{\top}\mathbf{h}+b_{2},\qquad
  a=\sigma(z_{2}),
\]
then $L$ is the binary cross-entropy of $a$ and $y$. Every link in the chain has its own local
derivative, and we already have them all:
\[
  \frac{\partial L}{\partial z_{2}}=a-y\ \text{(previous proposition)},\qquad
  \frac{\partial z_{2}}{\partial h_{i}}=w_{2}[i],
\]
\[
  \frac{\partial h_{i}}{\partial z_{1}[i]}=\mathbb{1}\bigl[z_{1}[i]>0\bigr],\qquad
  \frac{\partial z_{1}[i]}{\partial W_{1}[i,j]}=x_{j}.
\]
Multiplying backward along that chain gives, one by one,
\[
  \bar{z}_{2}=a-y,\qquad
  \nabla_{\mathbf{w}_{2}}L=\bar{z}_{2}\,\mathbf{h},\qquad
  \frac{\partial L}{\partial b_{2}}=\bar{z}_{2},
\]
then the signal that reaches the hidden layer
\[
  \bar{\mathbf{h}}=\bar{z}_{2}\,\mathbf{w}_{2},
  \qquad
  \bar{\mathbf{z}}_{1}=\bar{\mathbf{h}}\odot\mathbb{1}\bigl[\mathbf{z}_{1}>0\bigr],
\]
with $\odot$ the componentwise product. Note that $\bar{\mathbf{z}}_{1}$ is $\bar{\mathbf{h}}$
shaved by the ReLU gate: components whose gate is closed become exactly zero, not small. After
that
\[
  \nabla_{\mathbf{W}_{1}}L=\bar{\mathbf{z}}_{1}\mathbf{x}^{\top},\qquad
  \nabla_{\mathbf{b}_{1}}L=\bar{\mathbf{z}}_{1},\qquad
  \nabla_{\mathbf{x}}L=\mathbf{W}_{1}^{\top}\bar{\mathbf{z}}_{1}.
\]
The outer-product form of $\nabla_{\mathbf{W}_{1}}L$ is no accident: the derivative
$\partial z_{1}[i]/\partial W_{1}[i,j]=x_{j}$ does not depend on $i$, so the whole $i$-th row is
$\bar{z}_{1}[i]$ times the vector $\mathbf{x}$. Every weight receives a gradient whose size is
proportional to the input it multiplies, and that is why large-scale features automatically
receive large steps.

None of those formulas depends on $d=4$ or on two hidden units; the pattern is the same for any
layer.

The cost is comparable: one forward step needs one matrix-vector product per layer, one backward
step one transposed product and one outer product per layer. A centered finite difference for $P$
parameters needs $2P$ forward steps, whereas backpropagation needs about two. For $P=13$ in this
chapter, twenty-six against two; for a million parameters, two million against two, and that is
the whole difference. Nielsen names this observation as the reason backpropagation is regarded
as a great discovery, not merely an algebraic trick.

\begin{tangan}
Take three sheets of paper and a pencil, and label them: INPUT, WEIGHTS, OUTPUT. On the INPUT
sheet write four numbers: $1$; $0.5$; $-1$; $2$. On the WEIGHTS sheet write the two weights to
the output: $0.8$ and $-0.6$. On the OUTPUT sheet write $0.7446$, the probability the network
gives, and remember that the true label is $1$.

Lay the three sheets in a row from left to right: INPUT, WEIGHTS, OUTPUT. Now work backward,
always from right to left. On the OUTPUT sheet, subtract $0.7446$ from $1$ and write $-0.2554$.
Move to the WEIGHTS sheet: multiply $-0.2554$ by $0.8$, write $-0.2043$; multiply by $-0.6$,
write $0.1532$. Move to the INPUT sheet: multiply $-0.2043$ by each number on that sheet, then
write the row: $-0.2043$; $-0.1022$; $0.2043$; $-0.4086$.

What you must see: every time a sheet moves one step to the left, the sign of the number can
flip, and its size follows the contents of the sheet you are holding. If you reverse the
direction, from INPUT to OUTPUT, the order of work is wrong and the result means nothing.
\end{tangan}


\section{Initialization and gradient descent}
\tingkat{intermediate}

Once the gradient exists, the next step is simple: move a little in the direction that lowers the
loss. For all parameters $\boldsymbol{\theta}$,
\[
  \boldsymbol{\theta}\leftarrow\boldsymbol{\theta}-\eta\,\nabla_{\boldsymbol{\theta}}L,
\]
with $\eta$ the learning rate. For $B$ examples and the average loss
$L=\frac{1}{B}\sum_{i}L_{i}$, the gradient is also the average of the per-example gradients; the
backpropagation formula does not change, only how many backward passes. The conditions and
bounds on $\eta$ are in Chapter~23; because the network is nonconvex, the guarantees there are
only a benchmark.

\kamus{stochastic gradient descent}{penurunan gradien stokastik}{Walking down a slope using a
slope estimate from a small part of the data, not from all of the data.}\index{gradient descent}

Three variants deserve mention. SGD takes a small $B$ and updates the parameters every time; the
gradient is noisy, and the spread of that noise shrinks like $1/\sqrt{B}$ because the variance of
the average of $B$ independent variables is $\sigma^{2}/B$. Momentum keeps a velocity,
\[
  \mathbf{v}\leftarrow\beta\mathbf{v}+\mathbf{g},\qquad
  \boldsymbol{\theta}\leftarrow\boldsymbol{\theta}-\eta\mathbf{v},
\]
so that consistent directions accumulate and oscillating directions cancel. Adam normalizes every
coordinate with an estimate of its second moment,
\[
  \mathbf{m}\leftarrow\beta_{1}\mathbf{m}+(1-\beta_{1})\mathbf{g},\qquad
  \mathbf{v}\leftarrow\beta_{2}\mathbf{v}+(1-\beta_{2})\mathbf{g}^{2},
\]
\[
  \hat{\mathbf{m}}=\frac{\mathbf{m}}{1-\beta_{1}^{t}},\qquad
  \hat{\mathbf{v}}=\frac{\mathbf{v}}{1-\beta_{2}^{t}},\qquad
  \boldsymbol{\theta}\leftarrow\boldsymbol{\theta}
  -\eta\,\frac{\hat{\mathbf{m}}}{\sqrt{\hat{\mathbf{v}}}+\epsilon},
\]
with $t$ the step number. The factors $1-\beta_{1}^{t}$ and $1-\beta_{2}^{t}$ are called bias
correction, and their purpose is to remove the tendency of the moments toward zero in the early
steps. On the first step, $\hat{m}_{1}=g$ and $\hat{v}_{1}=g^{2}$, so the update is exactly
$-\eta\,g/(\lvert g\rvert+\epsilon)$: its size is almost exactly $\eta$, however large the
gradient. That property makes Adam scale-insensitive, and at the same time explains its weakness:
because its step does not shrink when the gradient shrinks, Adam is not guaranteed to converge on
all convex problems, and the counterexample can be read in Reddi, Kale, and Kumar (2018).

What remains, and what is most often underestimated, is the starting point.

\begin{proposisi}[Symmetry that never breaks]
If two rows of $\mathbf{W}_{1}$ are identical and $b_{1}[1]=b_{1}[2]$ at the start of training,
then both stay identical after every gradient update, for every data set and every learning
rate.\index{initialization}
\end{proposisi}

\begin{proof}
We prove it by induction on the update step. Suppose the first and second rows of
$\mathbf{W}_{1}$ are the same and $b_{1}[1]=b_{1}[2]$. Then for every input $z_{1}[1]=z_{1}[2]$,
so both ReLU gates make the same decision and $\bar{z}_{1}[1]=\bar{z}_{1}[2]$. Because
$\nabla_{\mathbf{W}_{1}}L=\bar{\mathbf{z}}_{1}\mathbf{x}^{\top}$, those two rows of the gradient
are identical, and so are the two components of $\nabla_{\mathbf{b}_{1}}L$. An update with the
same rate preserves the equality. For a set of examples, the total gradient is the sum of the
per-example gradients, so the equality survives that too.
\end{proof}

The consequence is blunt: eight units initialized uniformly behave like one unit, however long
it is trained.

Zero initialization is worse still. If all of $\mathbf{W}_{1}$ and $\mathbf{b}_{1}$ are zero,
then $\mathbf{z}_{1}=\mathbf{0}$, and because the gate is $\mathbb{1}[z>0]$, all gates are
closed: the gradient with respect to $\mathbf{W}_{1}$ is exactly zero, and the weights stay zero
forever. The only thing that can still learn is $b_{2}$, and the result is a constant prediction
equal to the class proportion. On the lab data, the final loss is $0.690923$, exactly the label
entropy. Eight units, zero benefit.

\begin{tantangan}
Design a 2--2--1 network with ReLU in the hidden layer whose gradient is exactly zero on all six
weights and both biases of the first layer from the first step, while $\mathbf{W}_{1}$ is not
zero, and the network is still wrong on at least one point. Only $b_{2}$ may move. Show why your
choice works.

Write your answer as a table: unit, $z_{1}[i]$ for each example, its gate state, and the gradient
that comes in.

Hint 1: what freezes the weights is not the value of the weights, but the ReLU gate closing for
all examples.

Hint 2: use data in the nonnegative quadrant, then take
$\mathbf{W}_{1}=-\mathbf{1}\mathbf{1}^{\top}$ and $\mathbf{b}_{1}=\mathbf{0}$. Then
$z_{1}[i]=-(x_{1}+x_{2})\le0$, so $\mathbf{h}=\mathbf{0}$,
$\bar{\mathbf{z}}_{1}=\mathbf{0}$, and $\nabla_{\mathbf{W}_{1}}L=\mathbf{0}$ even though
$\mathbf{W}_{1}\neq\mathbf{0}$.

If you get stuck, go back to the section Initialization and gradient descent.
\end{tantangan}


Random initialization is therefore not decoration, but a condition for those units to explore
different functions. The next question: how large. The answer comes from computing variances. If
the components of $\mathbf{x}$ are independent with mean zero and variance $\sigma_{x}^{2}$, and
the weights are independent with variance $\sigma_{w}^{2}$, then
\[
  \operatorname{Var}\bigl(z_{1}[i]\bigr)
  =\sum_{j=1}^{n_{\mathrm{in}}}\operatorname{Var}\bigl(W_{1}[i,j]x_{j}\bigr)
  =n_{\mathrm{in}}\,\sigma_{w}^{2}\sigma_{x}^{2}.
\]
In order for that variance not to depend on $n_{\mathrm{in}}$, we need
$\sigma_{w}^{2}\propto1/n_{\mathrm{in}}$. The He initialization uses
$\sigma_{w}^{2}=2/n_{\mathrm{in}}$, and the number $2$ is not decoration: for symmetric $z$,
ReLU cuts away half the distribution, so $\mathbb{E}[h^{2}]=\operatorname{Var}(z)/2=\sigma_{x}^{2}$,
and that variance survives layer after layer. The Xavier initialization uses
$\sigma_{w}^{2}=2/(n_{\mathrm{in}}+n_{\mathrm{out}})$, a compromise between preserving the forward
signal and the backward signal. That compromise is exact only when
$n_{\mathrm{in}}=n_{\mathrm{out}}$.

For sigmoid layers, an initialization that is too large punishes immediately: on a two-layer
network with standard deviation $1.0$, the initial loss reaches $2.223899$, whereas with standard
deviation $0.5$ it is only $1.220409$. The reverse is more serious: with standard deviation $0.5$
and eight layers, training stops at accuracy $0.5333$. Initialization is not an administrative
detail; it determines whether the signal is still there when it reaches the first layer.

\section{Worked example}
\tingkat{intermediate}

\begin{resep}
\begin{enumerate}
  \item Start from the end: $\bar{z}_{2}=a-y$ for sigmoid with cross-entropy, or
        $\mathbf{p}-\mathbf{y}$ for softmax. Without this number there is nothing to divide.
  \item Output weights: $\nabla_{\mathbf{w}_{2}}L=\bar{z}_{2}\mathbf{h}$ and
        $\partial L/\partial b_{2}=\bar{z}_{2}$.
  \item Signal to the previous layer: $\bar{\mathbf{h}}=\bar{z}_{2}\mathbf{w}_{2}$, then shave
        with the gate
        $\bar{\mathbf{z}}_{1}=\bar{\mathbf{h}}\odot\mathbb{1}[\mathbf{z}_{1}>0]$.
  \item First-layer weights: $\nabla_{\mathbf{W}_{1}}L=\bar{\mathbf{z}}_{1}\mathbf{x}^{\top}$
        and $\nabla_{\mathbf{b}_{1}}L=\bar{\mathbf{z}}_{1}$; if needed, the input gradient
        $\nabla_{\mathbf{x}}L=\mathbf{W}_{1}^{\top}\bar{\mathbf{z}}_{1}$.
\end{enumerate}
A warning about order: step 3 must not come before step 1. $\bar{\mathbf{z}}_{1}$ is computed
from $\bar{z}_{2}$, and the gate uses $\mathbf{z}_{1}$ from the forward computation that has
already been stored. Swapping the two steps kills the entire gradient of the first-layer weights.
\end{resep}


We use one example and compute everything with exact numbers. The input is
$\mathbf{x}=(1.0;\,0.5;\,-1.0;\,2.0)$, the label is $y=1$, and the initial parameters are
\[
  \mathbf{W}_{1}=\begin{pmatrix}
  0.5&-0.5&0.25&1.0\\
  0.1&0.2&-0.3&0.4\end{pmatrix},\quad
  \mathbf{b}_{1}=\begin{pmatrix}0.1\\-0.2\end{pmatrix},\quad
  \mathbf{w}_{2}=\begin{pmatrix}0.8\\-0.6\end{pmatrix},\quad b_{2}=0.05 .
\]
That is all the parameters.

\emph{Forward step.} The first unit receives
$0.5(1)-0.5(0.5)+0.25(-1)+1(2)+0.1=0.5-0.25-0.25+2+0.1=2.1$.
The second unit receives
$0.1(1)+0.2(0.5)-0.3(-1)+0.4(2)-0.2=0.1+0.1+0.3+0.8-0.2=1.1$.
Both are positive, so $\mathbf{h}=(2.1;\,1.1)$ without clipping. The output score is
$z_{2}=0.8(2.1)-0.6(1.1)+0.05=1.68-0.66+0.05=1.07$, so
$a=\sigma(1.07)=0.7445969158$ and the loss is $L=-\ln a=0.2949122596$.

\emph{Backward step.} The output adjoint is $\bar{z}_{2}=a-y=-0.2554030842$. From there
\[
  \nabla_{\mathbf{w}_{2}}L=\bar{z}_{2}\mathbf{h}
  =(-0.5363464768;\,-0.2809433926),\qquad
  \frac{\partial L}{\partial b_{2}}=-0.2554030842 .
\]
The signal to the hidden layer is
$\bar{\mathbf{h}}=\bar{z}_{2}\mathbf{w}_{2}=(-0.2043224673;\,0.1532418505)$; both gates are
open, so $\bar{\mathbf{z}}_{1}=\bar{\mathbf{h}}$. The outer product gives
\[
  \nabla_{\mathbf{W}_{1}}L=
  \begin{pmatrix}
  -0.2043224673&-0.1021612337&0.2043224673&-0.4086449347\\
  0.1532418505&0.0766209253&-0.1532418505&0.3064837010
  \end{pmatrix},
\]
and $\nabla_{\mathbf{b}_{1}}L=(-0.2043224673;\,0.1532418505)$. Two quick checks: the second row
is exactly $-0.75$ times the first, because $-0.6/0.8=-0.75$; and every column is proportional
to $x_{j}$, so the fourth column is twice the first, the second is half of it, and the third is
its negative. The Euclidean norm of the whole gradient is $0.9796964023$.

One gradient-descent step with $\eta=0.5$ updates every parameter according to
$\theta\leftarrow\theta-\eta\,\partial L/\partial\theta$, giving
\[
  \mathbf{W}_{1}'=\begin{pmatrix}
  0.6021612337&-0.4489193832&0.1478387663&1.2043224673\\
  0.0233790747&0.1616895374&-0.2233790747&0.2467581495
  \end{pmatrix},
\]
\[
  \mathbf{b}_{1}'=(0.2021612337;\,-0.2766209253),\qquad
  \mathbf{w}_{2}'=(1.0681732384;\,-0.4595283037),
\]
\[
  b_{2}'=0.1777015421 .
\]
After that update $\mathbf{z}_{1}'=(2.8406689;\,0.5444983)$ and $a'=0.9508189702$, so the loss
falls to $0.0504315919$; the difference is $-0.2444806677$, or about $83\%$ of the initial loss.

The number $83\%$ invites a misreading. What is measured is a single example, with no test data,
on a landscape that is not convex. The valid conclusion is not ``the rate $0.5$ is good'', but
``this gradient direction lowers the loss on this example''.

\kertas{Check one row of the lab table by hand. Take $W_{1}[2,1]$ and compute
$\bar{z}_{1}[2]\,x_{1}$ to get $0.1532418505$. Then compute the centered finite difference for
that parameter yourself: run one forward step with $W_{1}[2,1]+10^{-6}$ and one forward step with
$W_{1}[2,1]-10^{-6}$, divide the difference of the losses by $2\times10^{-6}$, and compare with
the analytic number above. After that, explain why the result is not exactly the same, and why
the sign of the difference can change if $h$ is made smaller, $10^{-12}$.}

\section{Computational lab}
\tingkat{advanced}

The three blocks are run with Python 3.14.4 and NumPy 2.5.2, without machine-learning libraries;
the output is copied as it stands.

The first block checks the hand computation above, then compares every analytic gradient with a
centered finite difference at $h=10^{-6}$.

\begin{lab}
\begin{verbatim}
import numpy as np
np.set_printoptions(precision=10, suppress=True, linewidth=76)

# 4-2-1 network: ReLU hidden, sigmoid output
x  = np.array([1.0, 0.5, -1.0, 2.0])
y  = 1.0
W1 = np.array([[0.5, -0.5, 0.25, 1.0],
               [0.1,  0.2, -0.3,  0.4]])
b1 = np.array([0.1, -0.2])
w2 = np.array([0.8, -0.6])
b2 = 0.05

def forward(x, W1, b1, w2, b2):
    z1 = W1 @ x + b1
    h  = np.maximum(z1, 0.0)
    o  = w2 @ h + b2
    p  = 1.0/(1.0+np.exp(-o))
    return z1, h, o, p, -(y*np.log(p) + (1-y)*np.log(1-p))

def backward(x, W1, b1, w2, b2):
    z1, h, o, p, L = forward(x, W1, b1, w2, b2)
    d   = p - y                    # dL/dz2 for sigmoid + cross-entropy
    dz1 = d*w2*(z1 > 0)            # ReLU gate
    return dict(W1=np.outer(dz1, x), b1=dz1, w2=d*h, b2=d,
                x=W1.T @ dz1), L, p, z1, h, dz1

g, L, p, z1, h, dz1 = backward(x, W1, b1, w2, b2)
print("z1 =", z1, "  h =", h)
print("p = %.12f  L = %.12f" % (p, L))
print("dz1 =", dz1)
print("g_w2 =", g["w2"], " g_b2 = %.12f" % g["b2"])
for i in range(2):
    print("g_W1[%d] =" % i, np.array2string(g["W1"][i], precision=10,
                                            suppress_small=False))
print("g_x =", g["x"])

# centered finite-difference check, one parameter at a time
eps = 1e-6
def loss_at(nama, idx, delta):
    W, B, V, c = W1.copy(), b1.copy(), w2.copy(), b2
    if nama == "W1":  W[idx] += delta
    if nama == "b1":  B[idx] += delta
    if nama == "w2":  V[idx] += delta
    if nama == "b2":  c += delta
    return forward(x, W, B, V, c)[4]

daftar = [("W1", i, j) for i in range(2) for j in range(4)]
daftar += [("b1", i) for i in range(2)] + [("w2", i) for i in range(2)]
print("\nparameter     analytic        numeric     relative error")
for nama, *idx in daftar + [("b2",)]:
    idx = tuple(idx)
    ana = {"W1": g["W1"], "b1": g["b1"], "w2": g["w2"], "b2": np.array(g["b2"])}
    a = ana[nama][idx] if idx else ana[nama]
    n = (loss_at(nama, idx, eps) - loss_at(nama, idx, -eps))/(2*eps)
    rel = abs(a-n)/max(abs(a), abs(n), 1e-300)
    print("%-10s %15.10f %15.10f %12.2e" %
          (nama + "".join("[%d]" % i for i in idx), a, n, rel))

# one gradient-descent step with rate 0.5
eta = 0.5
W1n, b1n = W1 - eta*g["W1"], b1 - eta*g["b1"]
w2n, b2n = w2 - eta*g["w2"], b2 - eta*g["b2"]
z1b, hb, ob, pb, Lb = forward(x, W1n, b1n, w2n, b2n)
print("\none step: p = %.12f  L = %.12f  difference = %.12f"
      % (pb, Lb, L - Lb))
for i in range(2):
    print("W1'[%d] =" % i, np.array2string(W1n[i], precision=10))
print("w2' =", w2n, " b2' = %.12f" % b2n)

# gradient with respect to the input: analytic vs finite difference
E = np.eye(4)
gx = np.array([(forward(x + eps*E[j], W1, b1, w2, b2)[4]
              - forward(x - eps*E[j], W1, b1, w2, b2)[4])/(2*eps)
               for j in range(4)])
print("\ng_x analytic    =", g["x"])
print("g_x finite diff =", gx)
print("relative error   =", ["%.1e" % v for v in np.abs(gx-g["x"])/np.abs(g["x"])])
\end{verbatim}
The actual output:
\begin{verbatim}
z1 = [2.1 1.1]   h = [2.1 1.1]
p = 0.744596915827  L = 0.294912259605
dz1 = [-0.2043224673  0.1532418505]
g_w2 = [-0.5363464768 -0.2809433926]  g_b2 = -0.255403084173
g_W1[0] = [-0.2043224673 -0.1021612337  0.2043224673 -0.4086449347]
g_W1[1] = [ 0.1532418505  0.0766209253 -0.1532418505  0.306483701 ]
g_x = [-0.0868370486  0.1328096038 -0.097053172  -0.1430257271]

parameter     analytic        numeric     relative error
W1[0][0]     -0.2043224673   -0.2043224674     1.17e-10
W1[0][1]     -0.1021612337   -0.1021612337     2.53e-10
W1[0][2]      0.2043224673    0.2043224674     1.17e-10
W1[0][3]     -0.4086449347   -0.4086449347     4.95e-11
W1[1][0]      0.1532418505    0.1532418505     7.22e-11
W1[1][1]      0.0766209253    0.0766209252     8.33e-10
W1[1][2]     -0.1532418505   -0.1532418505     7.22e-11
W1[1][3]      0.3064837010    0.3064837010     1.09e-10
b1[0]        -0.2043224673   -0.2043224674     1.17e-10
b1[1]         0.1532418505    0.1532418506     6.16e-10
w2[0]        -0.5363464768   -0.5363464768     9.80e-11
w2[1]        -0.2809433926   -0.2809433925     2.41e-10
b2           -0.2554030842   -0.2554030842     3.59e-11

one step: p = 0.950818970185  L = 0.050431591881  difference = 0.244480667724
W1'[0] = [ 0.6021612337 -0.4489193832  0.1478387663  1.2043224673]
W1'[1] = [ 0.0233790747  0.1616895374 -0.2233790747  0.2467581495]
w2' = [ 1.0681732384 -0.4595283037]  b2' = 0.177701542086

g_x analytic    = [-0.0868370486  0.1328096038 -0.097053172  -0.1430257271]
g_x finite diff = [-0.0868370486  0.1328096037 -0.097053172  -0.1430257272]
relative error   = ['7.7e-10', '5.0e-10', '1.4e-10', '9.8e-11']
\end{verbatim}
\end{lab}

That table is the proof that was asked for. All thirteen parameters agree to nine or ten digits,
with the largest relative error $8.33\times10^{-10}$ at $W_{1}[1,2]$. That number should not be
zero: the centered finite difference carries a truncation error of order $h^{2}$ and a rounding
error of order $\varepsilon_{\mathrm{mach}}/h$ (Chapter 22), so there is a best $h$ that cannot
make the error zero. For $h=10^{-6}$ an order of $10^{-10}$ is exactly what is expected.
 

The gradient with respect to the input agrees to $7.7\times10^{-10}$ on the first component. That
number means: raising the second feature slightly raises the loss fastest, whereas raising the
fourth feature lowers it fastest. In Chapter 29 this direction is used in reverse, to raise the
loss as fast as possible on inputs the attacker is allowed to change.

The second block trains a 2--8--1 network from scratch on two-dimensional data with an XOR
pattern, then compares it with a linear baseline and with two defective initializations.

\begin{lab}
\begin{verbatim}
import numpy as np

rng = np.random.default_rng(7)
n = 240
X = rng.uniform(-1.0, 1.0, size=(n, 2))
y = (X[:, 0]*X[:, 1] > 0).astype(float)        # label 1 when the signs agree
print("examples =", n, " class 1 =", int(y.sum()))

def sig(o):
    return 1.0/(1.0+np.exp(-o))

def bce(p, t):
    p = np.clip(p, 1e-12, 1-1e-12)
    return float(-np.mean(t*np.log(p) + (1-t)*np.log(1-p)))

def init(mode, rs, nh=8):
    if mode == "nol":
        return np.zeros((nh, 2)), np.zeros(nh), np.zeros(nh), 0.0
    if mode == "seragam":
        return np.full((nh, 2), 0.5), np.zeros(nh), np.full(nh, 0.5), 0.0
    return (rs.normal(0, np.sqrt(2/2), (nh, 2)), np.zeros(nh),
            rs.normal(0, np.sqrt(2/nh), nh), 0.0)          # He

def fwd(X, W1, b1, w2, b2):
    Z1 = X @ W1.T + b1
    H  = np.maximum(Z1, 0.0)
    return Z1, H, sig(H @ w2 + b2)

def grads(X, t, W1, b1, w2, b2):
    Z1, H, P = fwd(X, W1, b1, w2, b2)
    d  = (P - t)/len(t)
    dZ = np.outer(d, w2)*(Z1 > 0)
    return dZ.T @ X, dZ.sum(0), H.T @ d, d.sum()

def latih(mode, langkah=3000, lr=0.5, seed=0):
    W1, b1, w2, b2 = init(mode, np.random.default_rng(seed))
    jejak = []
    for t in range(langkah + 1):
        gW1, gb1, gw2, gb2 = grads(X, y, W1, b1, w2, b2)
        W1, b1 = W1 - lr*gW1, b1 - lr*gb1
        w2, b2 = w2 - lr*gw2, b2 - lr*gb2
        if t % 500 == 0:
            jejak.append((t, bce(fwd(X, W1, b1, w2, b2)[2], y)))
    Z1, H, P = fwd(X, W1, b1, w2, b2)
    mati = int(np.sum(np.all(H <= 0, axis=0)))
    return jejak, float(np.mean((P > 0.5) == (y > 0.5))), bce(P, y), mati

jejak, akur, L, mati = latih("he")
for t, v in jejak:
    print("step %4d   L = %.6f" % (t, v))
print("He      : accuracy = %.4f  dead units = %d/8" % (akur, mati))
for mode in ("nol", "seragam"):
    j, a, l, m = latih(mode)
    print("%-8s: L final = %.6f  accuracy = %.4f  dead units = %d/8"
          % (mode, l, a, m))

w, b = np.zeros(2), 0.0             # garis dasar linear, tanpa lapisan
for t in range(4000):
    p = sig(X @ w + b)
    w -= 0.5*(X.T @ (p - y))/n
    b -= 0.5*np.mean(p - y)
pl = sig(X @ w + b)
print("linear  : L = %.6f  accuracy = %.4f"
      % (bce(pl, y), float(np.mean((pl > 0.5) == (y > 0.5)))))
\end{verbatim}
The actual output:
\begin{verbatim}
examples = 240  class 1 = 112
step    0   L = 0.648815
step  500   L = 0.060245
step 1000   L = 0.036445
step 1500   L = 0.027164
step 2000   L = 0.022039
step 2500   L = 0.018717
step 3000   L = 0.016349
He      : accuracy = 1.0000  dead units = 0/8
nol     : L final = 0.690923  accuracy = 0.5333  dead units = 8/8
seragam : L final = 0.539475  accuracy = 0.7167  dead units = 0/8
linear  : L = 0.686130  accuracy = 0.7208
\end{verbatim}
\end{lab}

Its loss curve decreases at every sampling, from $0.648815$ to $0.016349$. Training accuracy
reaches $1.0000$ with eight units and zero dead units. The linear baseline, which uses the same
optimization and the same data but no hidden layer, stops at $0.7208$. That difference of $0.28$
is the price of a single separating plane, and it agrees with the XOR proposition: this pattern
is not linearly separable.

Figure~\ref{fig:kurva-kerugian} plots those seven points.

\begin{figure}[htbp]
\centering
\begin{tikzpicture}
\begin{axis}[
  width=12.0cm, height=5.6cm,
  xlabel={training step}, ylabel={loss $L$},
  xmin=-80, xmax=3080, ymin=0, ymax=0.75,
  axis lines=left,
  grid=major, grid style={layergrey!15},
  tick label style={font=\scriptsize},
  ytick={0,0.2,0.4,0.6}, yticklabels={$0$,$0.2$,$0.4$,$0.6$},
  label style={font=\small},
  legend style={font=\scriptsize, draw=none, fill=none,
                at={(0.98,0.98)}, anchor=north east},
]
\addplot[thick, accent, mark=*, mark size=1.7pt] coordinates {
  (0,0.648815) (500,0.060245) (1000,0.036445) (1500,0.027164)
  (2000,0.022039) (2500,0.018717) (3000,0.016349) };
\addplot[dashed, layergrey, domain=0:3000, samples=2] {0.693147};
\legend{training loss, $\ln 2=0.693147$}
\end{axis}
\end{tikzpicture}
\caption{The training loss curve of the 2--8--1 network in the second block, from $0.648815$ at
step $0$ to $0.016349$ at step $3000$. The dashed line is $\ln 2=0.693147$, the loss of a
uniform guess; the first point is already below it because it is measured after the first
update.}
\label{fig:kurva-kerugian}
\end{figure}


Zero initialization teaches a harsher lesson. Its final loss is $0.690923$, equal to the label
entropy to six decimals, and its accuracy is $0.5333$, namely the proportion of the majority
class; all eight units are dead. Uniform initialization gives $0.7167$, practically the same as
the linear baseline.

The third block measures two things: how many layers a sigmoid network can still learn through,
and how fast the signal shrinks from the output to the first layer.

\begin{lab}
\begin{verbatim}
import numpy as np

rng = np.random.default_rng(7)
n = 240
X = rng.uniform(-1.0, 1.0, size=(n, 2))
y = (X[:, 0]*X[:, 1] > 0).astype(float)

def sig(o):
    return 1.0/(1.0+np.exp(-o))

def bce(p, t):
    p = np.clip(np.ravel(p), 1e-12, 1-1e-12)
    t = np.ravel(t)
    return float(-np.mean(t*np.log(p) + (1-t)*np.log(1-p)))

def bangun(dalam, std, seed=0):
    rs = np.random.default_rng(seed)
    ukuran = [2] + [16]*(dalam-1) + [1]
    W = [rs.normal(0, std, (ukuran[i+1], ukuran[i])) for i in range(dalam)]
    return W, [np.zeros(ukuran[i+1]) for i in range(dalam)]

def maju(W, b, X):
    A = [X]
    for Wl, bl in zip(W[:-1], b[:-1]):
        A.append(sig(A[-1] @ Wl.T + bl))
    return A, sig(A[-1] @ W[-1].T + b[-1])

def mundur(W, b, X, t):
    A, P = maju(W, b, X)
    d = (P - t[:, None])/len(t)
    gW, gb = [None]*len(W), [None]*len(b)
    gW[-1], gb[-1] = d.T @ A[-1], d.sum(0)
    delta = d @ W[-1]
    for l in range(len(W)-1, 0, -1):
        delta = delta*A[l]*(1 - A[l])
        gW[l-1], gb[l-1] = delta.T @ A[l-1], delta.sum(0)
        if l > 1:
            delta = delta @ W[l-1]
    return gW, gb, bce(P, t), P

print("depth  std   L(start)   L(final)  accuracy  ||gW1||/||gWlast||")
for dalam in (1, 2, 4, 8):
    for std in (0.5, 1.0):
        W, b = bangun(dalam, std)
        gW, gb, L0, P = mundur(W, b, X, y)
        rasio = np.linalg.norm(gW[0])/np.linalg.norm(gW[-1])
        for t in range(2000):
            gW, gb, L, P = mundur(W, b, X, y)
            for l in range(dalam):
                W[l] -= 0.5*gW[l]
                b[l] -= 0.5*gb[l]
        akur = float(np.mean((np.ravel(P) > 0.5) == (y > 0.5)))
        print("%5d %5.1f %9.6f %9.6f %8.4f %14.2e"
              % (dalam, std, L0, bce(P, y), akur, rasio))
\end{verbatim}
The actual output:
\begin{verbatim}
depth  std   L(start)   L(final)  accuracy  ||gW1||/||gWlast||
    1   0.5  0.693190  0.686130   0.7208       1.00e+00
    1   1.0  0.693951  0.686130   0.7208       1.00e+00
    2   0.5  1.220409  0.332913   0.8417       1.66e-02
    2   1.0  2.223899  0.299519   0.8542       2.16e-02
    4   0.5  0.691597  0.147887   0.9458       5.77e-02
    4   1.0  0.701977  0.023491   0.9958       2.05e-01
    8   0.5  0.739920  0.690867   0.5333       1.26e-04
    8   1.0  0.875314  0.002383   1.0000       3.76e-03
\end{verbatim}
\end{lab}

The first three rows of that table are a promise and a warning at once. Depth one gives $0.7208$,
the same as the linear baseline in the previous block; depth two gives $0.8417$ and depth four
$0.9458$. But the last row with standard deviation $0.5$ gives $0.5333$ with loss $0.690867$:
eight sigmoid layers with a small initialization practically do not learn. The gradient-norm
ratio in the last column explains why: $1.26\times10^{-4}$, four orders of magnitude smaller than
at depth four. The same row with standard deviation $1.0$ succeeds perfectly. So the failure is
not about depth alone, but about the initial scale.

Two limits of these lab results must be stated. The data is only 240 self-made two-dimensional
examples, with no test set, so all of it is training accuracy; generalization is the business of
Chapter 26 and Chapter 27. All experiments use a single seed, so variation between seeds is
invisible. What may be concluded is only what was checked directly: backpropagation agrees with
finite differences, the loss decreases, zero and uniform initialization fail in ways that can be
explained, and a deep network with a small initialization loses its signal.

\begin{untukdosen}
The order that has worked: start from one separating line and four XOR points, then show the
2--2--1 construction on the board, and only after that the general definition. The proof that XOR
is impossible is usually accepted easily; what stops students is the backward step, so set aside
the most time for backpropagation, not for activation functions.

Two typical misconceptions. First, ``backpropagation is another algorithm, not the chain rule'',
so that students memorize formulas without being able to derive them for another architecture;
the cure is to ask them to re-derive the 4--2--1 network with three hidden units on the board.
Second, ``a relative error of $10^{-10}$ means the implementation is proved correct''; a small
error only checks agreement at one point, for one set of parameters, and does not check the sign,
the scaling, or the case of a closed gate.

The most helpful board activity: give a network with one hidden unit whose $z_{1}<0$, ask the
class to compute the gradient of the whole row $\mathbf{W}_{1}$, then ask what happens to that
unit after a hundred steps. The answer ``nothing happens, ever'' usually opens the discussion
about unit death and about why a large learning rate is dangerous.

Sign of readiness to go on: students can write the backpropagation for a network with two hidden
layers without looking at their notes, and can explain why zero initialization fails not because
of ``too little data'' but because of symmetry.
\end{untukdosen}

\section{Summary}

\begin{ringkasan}
\begin{itemize}
  \item A linear threshold can only separate sets that are linearly separable, and the four XOR
        points are not one of them. The proof does not depend on the algorithm, so adding data
        does not help. One hidden layer with ReLU is already enough for that pattern.
  \item Without a nonlinear activation, a layered network collapses into a single matrix. Sigmoid
        and tanh shrink the signal because their derivatives are at most $\tfrac14$ and $1$;
        through $L$ sigmoid layers the remaining signal is at most $4^{-L}$, namely
        $9.54\times10^{-7}$ for $L=10$. ReLU does not shrink, but it kills units whose input is
        always negative.
  \item Cross-entropy is paired with sigmoid or softmax so that the output adjoint is exactly
        $a-y$ or $\mathbf{p}-\mathbf{y}$. With the squared loss, the factor $a(1-a)$ is left
        behind and the signal shrinks exactly when the model is most confidently wrong; at $z=10$
        the ratio is about $2.2\times10^{4}$.
  \item Backpropagation is the chain rule of Chapter 12 run on the computational graph in reverse
        order. For 4--2--1, all the gradients are:
        $\bar{z}_{2}=a-y$, $\nabla_{\mathbf{w}_{2}}L=\bar{z}_{2}\mathbf{h}$,
        $\bar{\mathbf{z}}_{1}=\bar{z}_{2}\mathbf{w}_{2}\odot\mathbb{1}[\mathbf{z}_{1}>0]$,
        $\nabla_{\mathbf{W}_{1}}L=\bar{\mathbf{z}}_{1}\mathbf{x}^{\top}$,
        $\nabla_{\mathbf{x}}L=\mathbf{W}_{1}^{\top}\bar{\mathbf{z}}_{1}$. Its cost is about two
        forward steps, not $2P$. On the 4--2--1 example, those thirteen gradients agree with
        centered finite differences to a relative error of $8.33\times10^{-10}$.
  \item Initialization determines whether the network can learn. Zero closes all gates and leaves
        a prediction equal to the class proportion; uniform preserves the symmetry so that the
        effective capacity is one unit; too small shrinks the signal before the first layer; too
        large saturates the sigmoid. He uses $\sigma_{w}^{2}=2/n_{\mathrm{in}}$, Xavier
        $2/(n_{\mathrm{in}}+n_{\mathrm{out}})$.
  \item Gradient descent, SGD, momentum, and Adam differ only in how they use the same gradient.
        Adam's first step has a size of almost exactly $\eta$ however large the gradient, and that
        property explains both its robustness and its failure on some problems.
  
\end{itemize}
\end{ringkasan}

\section{Exercises}

\begin{latihan}
\begin{enumerate}
  \item For data with margin $\gamma=0.2$ and $\lVert\mathbf{x}_{i}\rVert\le1$, what is the upper
        bound on the number of perceptron updates according to the convergence theorem, and what
        is it for $\gamma=0.05$? Explain why that bound does not depend on the number of examples
        and why it is useless when $\gamma$ is small.

  \item Prove that a two-layer network with identity activation in the hidden layer is no stronger
        than a linear threshold, that is, that
        $\mathbf{W}_{2}\operatorname{id}(\mathbf{W}_{1}\mathbf{x}+\mathbf{b}_{1})+\mathbf{b}_{2}$
        can be written as a single linear threshold. Explain the consequence for further identity
        layers.

  

  \item Compute $\nabla_{\mathbf{W}_{1}}L$ for the network in the worked example if
        $\mathbf{x}=(0.5;\,-1.0;\,0.25;\,2.0)$ and $y=0$, with the same parameters. Determine
        which unit's gate is closed, and why that row has zero gradient without meaning that the
        loss is minimal.

  \item Assess this claim: \emph{because the largest relative error is only
        $8.33\times10^{-10}$, the backpropagation implementation in this chapter is proved
        correct.} Decide true, false, or not enough evidence, and name the additional checks
        required.

  \item Assess this claim: \emph{the universal approximation theorem guarantees that a network
        with one hidden layer can reach an arbitrarily small error on test data.} Decide true,
        false, or not enough evidence, and name two things the theorem does not state.

  \item With $\bar{z}_{2}=-0.2554030842$, compute one Adam step for $b_{2}=0.05$ with
        $\eta=0.01$, $\beta_{1}=0.9$, $\beta_{2}=0.999$, $\epsilon=10^{-8}$, at $t=1$. Show that
        the step size is practically $\eta$, then compute the same step if the gradient is $100$
        times larger.

  \item Show that for two ReLU units in $\mathbb{R}^{2}$ the output is a piecewise linear
        function, and that the input plane is divided into at most four regions with a linear
        formula each. How many regions for $k$ ReLU units in $\mathbb{R}^{2}$ in general
        position? (Hint: each unit gives one line $z_{1}[i]=0$, and inside each cell all gates
        stay fixed.)

  \item Compute the gradient with respect to the input for the worked example, then determine the
        perturbation $\boldsymbol{\delta}$ with $\lVert\boldsymbol{\delta}\rVert_{\infty}\le0.05$
        that raises the loss most according to the linearization, namely
        $\boldsymbol{\delta}=0.05\operatorname{sign}(\nabla_{\mathbf{x}}L)$. Compute the loss at
        $\mathbf{x}+\boldsymbol{\delta}$ and compare it with the linear prediction
        $L+\nabla_{\mathbf{x}}L^{\top}\boldsymbol{\delta}$.

  \item Without running a program, predict what happens if all of $\mathbf{b}_{1}$ is initialized
        to $+10$ while $\mathbf{W}_{1}$ stays small and random, using the saturation table; then
        if $\mathbf{b}_{1}=-10$. Check your prediction by running the second lab block after
        changing one line.
\end{enumerate}
\end{latihan}

Hints for the three hardest items. Item 2: write $\mathbf{W}=\mathbf{W}_{2}\mathbf{W}_{1}$ and
$\mathbf{b}=\mathbf{W}_{2}\mathbf{b}_{1}+\mathbf{b}_{2}$, then compare with the definition of a
linear threshold. Item 7: compute the number of regions for $k=1,2,3$ (they are $2,4,7$), then look for the
pattern $1+k+k(k-1)/2$. Item 9: the second lab block uses ReLU, so $\mathbf{b}_{1}=+10$ saturates nothing; what
saturates is $\mathbf{b}_{1}=-10$, when every unit is dead and its derivative zero. In the
third block, which uses the sigmoid, $z\approx10$ makes the derivative of the hidden layer
about $4.5\times10^{-5}$, so the gradient to $\mathbf{W}_{1}$ shrinks proportionally, while
$\mathbf{b}_{1}$ still receives the full gradient.

\begin{untukpeneliti}
What is really proved in this chapter: the impossibility of linear separation for XOR; the bound
$\lfloor1/\gamma^{2}\rfloor$ on the perceptron, with a complete proof; the derivative of the
sigmoid and the bound $\tfrac14$; the equality $\partial L/\partial z=a-y$ for sigmoid with
cross-entropy, and $\mathbf{p}-\mathbf{y}$ for softmax, both with proofs; the adjoint recursion as
the chain rule; and the proposition on initialization symmetry. What is borrowed without proof:
the universal approximation theorem, with the original proofs in Cybenko (1989) and Hornik
(1991), both using the representation theorem for continuous linear functionals. What is
popularized beyond its proof: ``a network can approximate anything'' as a practical promise, when
it is a statement about the density of functions, not about the reachability of parameters by
optimization; and ``Adam is always better'', when the counterexample is in Reddi, Kale, and Kumar
(2018).

The primary readings that shaped this chapter: Nielsen sections 2.4 and 2.8 (the four
backpropagation equations and why it is cheap) and section 3.1 (cross-entropy, softmax);
Goodfellow et al. section 6.4.1 (universal approximation) and 6.5 (backpropagation). For
initialization, Glorot and Bengio (2010) and He et al. (2015); for Adam, Kingma and Ba (2015).

Limits of validity: all computations in this chapter are for one example on a network with one
hidden layer, with the convention that the ReLU derivative is zero at the break point. There is no
statement about generalization here; that is the content of Chapter 26. Open questions:
why wide networks are easier to train than narrow ones at the same error, why a network with far
more parameters than examples still generalizes, and whether there is an initialization that is
optimal for any depth.
\end{untukpeneliti}

\peta{This chapter follows Nielsen, \emph{Neural Networks and Deep Learning}: sections 1.1
(perceptron), 1.2 (sigmoid neuron), 1.5 (gradient descent), 2.1 to 2.5 (matrix notation, the two
assumptions about the loss function, the Hadamard product, the four backpropagation equations
with their proofs), 2.8 (why backpropagation is fast), 3.1 (cross-entropy, softmax), and 3.3
(weight initialization). For the general formulation and for the side Nielsen does not prove,
Goodfellow, Bengio, and Courville, \emph{Deep Learning} is used: section 6.2 (gradient-based
learning, loss functions, and output units), 6.3 (hidden units and the choice of activation),
6.4.1 (universal approximation and depth), 6.5 (backpropagation and computational graphs), 8.3
(basic algorithms: SGD, momentum), 8.4 (parameter initialization strategies), and 8.5 (algorithms
with adaptive learning rates, including Adam). For the gradient with respect to the input as a
basis for attacks, see the discussion of FGSM in Chapter 27.}
